\section{The effective policy includes the harness}\label{sec:C20:policy}
Let $H_t$ denote the recorded interaction history, $C_t$ the constructed model context, $M_t$ persistent memory, and $E_t$ environment state. A harness configuration $h$ includes a context constructor $g_h$, action parser $p_h$, scheduler, authorization checks, and stopping policy. Let $S_t$ be the recorded execution state, defined in the next section. Let $\mathcal Z$ be the countable set of stopped output records and $\mathcal A$ the dispatch-action space augmented with a no-effect rejection outcome. Assume termination for this single-call model; otherwise nontermination needs its own mass as in Chapter~\ref{ch:C02}. For $z_t\in\mathcal Z$ and $a_t\in\mathcal A$, the induced distribution is
\begin{equation}
 \overline\Pi_{\theta,\delta,h}(a_t\mid S_t)=\sum_{z\in\mathcal Z}
 Q_{\theta,\delta}(z\mid C_t)
 K_h(a_t\mid p_h(z),S_t;\alpha_t).
\end{equation}
Here $Q_{\theta,\delta}$ is the normalized stopped law of the actual decoder, $C_t=g_h(S_t)$ uses only the declared policy-visible projection, and $K_h$ is the parser/enforcement kernel under current authority epoch $\alpha_t$. Its output is a dispatched action or a no-effect outcome; dispatch is not remote commitment. It can reject malformed output or substitute a deterministic fallback. A request for model repair is instead a transition to another inference call, not a free operation inside this single-call expression. Its likelihood, latency and cost belong to the subsequent trajectory. This distinction prevents a repair-enabled harness from being analyzed as though it executed only the final valid response. A parser that silently normalizes an ambiguous filename makes a policy decision; it is not merely formatting infrastructure.

This decomposition explains why two products using the same checkpoint can differ substantially. They may expose different tools, truncate different observations, retry at different boundaries, or terminate after different evidence. The code-as-harness survey organizes these mechanisms through interfaces, execution feedback, memory, and coordination~\cite{corpus260518747}. Its taxonomy is useful architectural synthesis, but does not establish that every additional mechanism improves performance. Additional machinery can increase failure opportunities and conceal the origin of a gain.

A clean causal comparison specifies the treatment as a configuration change. Suppose the treatment adds a planning file and a summarization tool. Its measured effect combines extra inference, altered context, persistent state, and a new action interface. Attributing the result to ``planning'' requires controls that preserve the other differences. One control allocates the same additional model calls to independent attempts; another retains the file interface but disables learned updates; a third supplies an oracle summary to estimate headroom. None is sufficient alone because the treatment can interact with task length and tool latency.

Harness adaptation creates a second policy. RRSI searches over harness changes while holding backbone weights fixed and constrains proposal and selection behavior~\cite{rrsi2026}. The important lesson is not a universal numerical gain, but that repeated selection against a finite evolve set can optimize quirks rather than transferable procedures. A runtime should therefore record the complete harness identity alongside the model and memory version. Reconstructing a trajectory with a later parser or prompt can change its interpretation even when every raw message is preserved.

For implementation, define the smallest state transition with an externally meaningful outcome. Avoid making an entire conversational turn the only transaction boundary: it can contain multiple tool effects and multiple opportunities for interruption. Also avoid recording every generated token as a committed action. Token streams are provisional inference products; an executable operation begins only after parsing, authorization, and allocation of a stable identity.

